% preface_section1_draft.tex — Section 1 of the Restored Preface
% To be integrated into preface.tex
% Macros assumed from main.tex: \cA, \barB, \gAmod, \MC, \cZ
\providecommand{\cA}{\mathcal{A}}
\providecommand{\barB}{\overline{B}}
\providecommand{\gAmod}{\mathfrak{g}_{\cA}^{\mathrm{mod}}}
\providecommand{\MC}{\text{MC}}

\section*{1.\quad The bar construction on algebraic curves}
% ====================================================================

\medskip

\noindent
The bar construction of a graded algebra is classical: it lives on the
cofree tensor coalgebra $T^c(s^{-1}\bar A)$ and its differential
encodes the multiplication.  Three extensions of this classical theory
are the subject of this monograph.  First, the multiplication of a
chiral algebra on a curve~$X$ is mediated by a logarithmic propagator
$d\log(z_i - z_j)$, and the bar differential becomes an integral
transform on the Fulton--MacPherson compactification of configuration
spaces.  Second, the bar coalgebra carries two coproducts: one
cocommutative (recording unordered collisions on~$\operatorname{Ran}(X)$,
producing the factorization coalgebra $\barB^{\Sigma}(\cA)$) and one
non-cocommutative (recording ordered collisions along
$\operatorname{Conf}_n^{<}(\mathbb R)$, producing the tensor coalgebra
$\barB^{\mathrm{ord}}(\cA)$).  The non-cocommutative coproduct is the
natural primitive, and its Koszul dual is a Yangian.  Third, on curves
of genus $g \ge 1$ the propagator deforms via the prime form, the bar
differential acquires curvature $\kappa(\cA)\cdot\omega_g$ from the
Hodge bundle, and the universal Maurer--Cartan element~$\Theta_{\cA}$
of the modular convolution algebra controls the genus tower.

% ====================================================================
\subsection*{The $1$-form}
% ====================================================================

\noindent
Let $X$ be a smooth algebraic curve over $\mathbf{C}$.  On the configuration space
$C_2(X) = X \times X \setminus \Delta$, define the logarithmic $1$-form
\begin{equation}\label{eq:pf1-eta}
\eta_{12} \;=\; d\log(z_1 - z_2) \;=\; \frac{dz_1 - dz_2}{z_1 - z_2}\,.
\end{equation}
This form has a simple pole along the diagonal $\Delta \subset X \times X$
with residue $1$.  The residue is intrinsic: it depends on no coordinate, no
metric, no level.  A double pole $dz/(z-w)^2$ transforms under coordinate
change and has no coordinate-independent residue.  A regular $1$-form on
$X \times X$ has residue zero along $\Delta$ and extracts nothing from the
collision.  The logarithmic derivative is the unique $1$-form on $C_2(X)$
with a first-order pole along $\Delta$ and conformally invariant residue.

On $C_3(X)$, the three pullbacks $\eta_{12}, \eta_{23}, \eta_{31}$ satisfy
the \emph{Arnold relation}
\begin{equation}\label{eq:pf1-arnold}
\eta_{12} \wedge \eta_{23} \;+\; \eta_{23} \wedge \eta_{31}
\;+\; \eta_{31} \wedge \eta_{12} \;=\; 0\,.
\end{equation}
The identity restates, in differential forms, the partial-fractions relation
\[
\frac{1}{(z_1 - z_2)(z_2 - z_3)}
\;+\; \frac{1}{(z_2 - z_3)(z_3 - z_1)}
\;+\; \frac{1}{(z_3 - z_1)(z_1 - z_2)}
\;=\; 0\,,
\]
which holds because the left side, viewed as a rational function of $z_1$
with $z_2, z_3$ fixed, has simple poles at $z_2$ and $z_3$ whose residues
sum to zero, hence vanishes.  Arnold proved that
\eqref{eq:pf1-arnold}, replicated across all triples
$\{i,j,k\} \subset \{1, \ldots, n\}$, generates the \emph{complete} ideal
of relations in $H^*(\mathrm{Conf}_n(\mathbf{C}), \mathbf{C})$.

Under a coordinate change $z \mapsto w(z)$, the propagator transforms as
$\eta_{ij} \mapsto \eta_{ij} + d\log\frac{w(z_i) - w(z_j)}{z_i - z_j}$.
Near the diagonal this becomes $d\log w'(z_i)$: the standard cocycle for
$\mathrm{Diff}(X)$.  The propagator is a BRST ghost for diffeomorphisms.


% ====================================================================
\subsection*{Chiral algebras on a curve}
% ====================================================================

Let $X$ be a smooth algebraic curve over~$\mathbb C$.  Write
$\mathcal D_X$ for the sheaf of differential operators on~$X$; a
$\mathcal D_X$-module is a sheaf on which $\mathcal D_X$ acts,
equivalently a sheaf with a flat connection.
A \emph{chiral algebra}~$\cA$ on~$X$ is a $\mathcal D_X$-module
equipped with a \emph{chiral bracket}
\[
\mu^{\mathrm{ch}}
\;\colon\;
j_*j^*(\cA\boxtimes\cA)
\;\longrightarrow\;
\Delta_!(\cA),
\]
where $j\colon X\times X\setminus\Delta
\hookrightarrow X\times X$ is the inclusion of the
complement of the diagonal and
$\Delta_!\cA$ is the $!$-pushforward along the diagonal
embedding $\Delta\colon X\hookrightarrow X\times X$
(the $\mathcal{D}$-module direct image; concretely,
$\Delta_!\cA \cong \cA \otimes \delta_\Delta$, sections
of~$\cA$ supported on the diagonal via the
$\delta$-function).  The source
$j_*j^*(\cA\boxtimes\cA)$ is the $\mathcal{D}$-module
of sections of the external product
$\cA \boxtimes \cA$ with poles of arbitrary order
along~$\Delta$; the chiral bracket extracts the
singular part and projects it onto the diagonal.

In local coordinates, the chiral bracket encodes the \emph{operator
product expansion}: for sections~$a,b\in\cA$,
\[
a(z)\,b(w)
\;\sim\;
\sum_{n\ge 0}\frac{a_{(n)}b(w)}{(z-w)^{n+1}}\,.
\]
The $n$-th products $a_{(n)}b\in\cA$ are the
\emph{Fourier modes} of the bracket, extracted by the
residue formula
$a_{(n)}b = \operatorname{Res}_{z=w}
[a(z)\,b(w)\,(z{-}w)^n]$.
The zeroth product $a_{(0)}b$ is the Lie-type
bracket (first-order pole); the first product $a_{(1)}b$
contributes the second-order pole, carrying the
symmetric/metric data.  All higher products ($n\ge 2$)
are determined by the $\mathcal D_X$-module structure,
i.e.\ by translation covariance.

The \emph{Borcherds identity} couples all pole orders into a single
algebraic constraint:
\[
\boxed{
\sum_{j\ge 0}\binom{m}{j}(a_{(n+j)}b)_{(m+k-j)}c
\;=\;
\sum_{j\ge 0}(-1)^j\binom{n}{j}
\Bigl(
a_{(m+n-j)}(b_{(k+j)}c)
-(-1)^n b_{(n+k-j)}(a_{(m+j)}c)
\Bigr).
}
\]
At $n=0$: the Jacobi identity for the Lie bracket~$a_{(0)}$.
At $m=0$: the commutator formula.  At $n=1$, $m=k=0$:
the derivation property of $a_{(0)}$ with respect to $a_{(1)}b$.
The Borcherds identity is the complete axiom system: locality,
Jacobi, and associativity in a single formula.


% ====================================================================
\subsection*{The bar complex on~$X$}
% ====================================================================

Place sections $a_1,\dots,a_n\in\cA$ at distinct points
$z_1,\dots,z_n\in X$.  Assign to each ordered pair $(i,j)$ the
logarithmic $1$-form
\[
\eta_{ij}
\;:=\;
d\log(z_i-z_j)
\;=\;
\frac{dz_i-dz_j}{z_i-z_j}\,.
\]
For a spanning tree~$T$ of the complete graph on $\{1,\dots,n\}$, with edge set~$E(T)$,
form the decorated tensor
\[
\omega_T(a_1,\dots,a_n)
\;:=\;
a_1(z_1)\otimes\cdots\otimes a_n(z_n)
\;\otimes\;
\bigwedge_{(i,j)\in E(T)}\eta_{ij}\,.
\]
The form $\bigwedge_{(i,j)\in E(T)}\eta_{ij}$ is a section of
$\Omega^{n-1}_{\log}$ on the configuration space
$\mathrm{Conf}_n(X)$; its degree is~$|E(T)|=n-1$ since~$T$ is a
tree.

The Fulton--MacPherson compactification
$\overline{\mathrm{FM}}_n(X)$ resolves $\mathrm{Conf}_n(X)$ by
iterated real oriented blowup along all diagonal strata.  This is a
\emph{blowup}, not the complement of the diagonal: points of
$\overline{\mathrm{FM}}_n(X)\setminus\mathrm{Conf}_n(X)$ record
the limiting direction and relative scales of approach as
points collide.  The boundary is a manifold with corners; the
codimension-one faces are indexed by subsets $S\subseteq\{1,\dots,n\}$
with $|S|\ge 2$, recording which points have collided.

The natural primitive is the \emph{ordered bar complex}
\[
\barB^{\mathrm{ord}}_X(\cA)
\;=\;
\bigoplus_{n\ge 1}
\bigl(s^{-1}\bar\cA\bigr)^{\otimes n}
\;\otimes\;
\Gamma\bigl(\overline{\mathrm{FM}}_n(X),\;
\Omega^{n-1}_{\log}\bigr),
\]
with no $\Sigma_n$-quotient.  Here $\bar\cA = \ker(\varepsilon\colon\cA\to\omega_X)$ is the
augmentation ideal (the kernel of the vacuum map), and
$s^{-1}$ denotes \emph{desuspension}: the shift
$|s^{-1}a| = |a| - 1$ in cohomological degree.  Throughout this
monograph, cohomological conventions are used: $|d| = +1$.

The ordered bar is a cofree tensor coalgebra: the \emph{deconcatenation coproduct}
\[
\Delta[s^{-1}a_1|\cdots|s^{-1}a_n]
\;=\;
\sum_{i=0}^{n}
[s^{-1}a_1|\cdots|s^{-1}a_i]
\;\otimes\;
[s^{-1}a_{i+1}|\cdots|s^{-1}a_n]
\]
splits an ordered sequence at every position: $n+1$ terms.  This is an
$E_1$-coalgebra: coassociative but \emph{not} cocommutative, reflecting the
linear ordering of the tensor factors.  The \emph{symmetric bar complex}
\[
\barB^{\Sigma}_X(\cA)
\;=\;
\bigoplus_{n\ge 1}
\bigl(s^{-1}\bar\cA\bigr)^{\otimes n}_{\Sigma_n}
\;\otimes\;
\Gamma\bigl(\overline{\mathrm{FM}}_n(X),\;
\Omega^{n-1}_{\log}\bigr)
\]
is the $\Sigma_n$-coinvariant quotient.  It carries the \emph{coshuffle coproduct},
summing over all $2^n$ bipartitions of an unordered set: coassociative and cocommutative,
an $E_\infty$-coalgebra, the factorization coalgebra of Beilinson and Drinfeld
on~$\operatorname{Ran}(X)$.

The bar differential $d_{\barB}$ decomposes into three pieces:
\[
d_{\barB}
\;=\;
d_{\mathrm{int}}+d_{\mathrm{res}}+d_{\mathrm{dR}}\,.
\]
The three components:
\begin{itemize}
\item $d_{\mathrm{int}}$: the internal differential of~$\cA$
(the translation operator, acting on each tensor factor).
\item $d_{\mathrm{res}}$: the \emph{residue differential}. For
a codimension-one boundary face $D_S$ of
$\overline{\mathrm{FM}}_n(X)$ corresponding to the collision
$\{z_i\}_{i\in S}\to z_0$, the residue along~$D_S$ extracts
the OPE coefficient: the iterated contraction
$a_{i_1(n_1)}(a_{i_2(n_2)}(\cdots))$ at the pole orders
dictated by the form $\bigwedge\eta_{ij}$.  In formulas: for
$S=\{i,j\}$,
\[
d_{\mathrm{res}}^{(ij)}
\bigl(
a_1\otimes\cdots\otimes a_n\otimes\eta_{ij}\wedge\omega'
\bigr)
\;=\;
\pm\,
a_1\otimes\cdots\otimes
\widehat{a_i}\otimes\cdots\otimes
\widehat{a_j}\otimes\cdots\otimes
a_{i(p)}a_j
\;\otimes\;
\omega',
\]
where $p$ is the pole order read from the form-theoretic
coefficient, and hats denote omission.
\item $d_{\mathrm{dR}}$: the de~Rham differential on the
logarithmic forms on~$\overline{\mathrm{FM}}_n(X)$.
\end{itemize}


% ====================================================================
\subsection*{The proof that $d_{\barB}^2=0$}
% ====================================================================

Expand $(d_{\mathrm{int}}+d_{\mathrm{res}}+d_{\mathrm{dR}})^2$
into nine terms.  The diagonal terms: $d_{\mathrm{int}}^2=0$
(since~$\cA$ is a cochain complex), $d_{\mathrm{dR}}^2=0$
(de~Rham), and $d_{\mathrm{res}}^2$ must be controlled.  The
cross-terms $d_{\mathrm{int}}\,d_{\mathrm{dR}}+
d_{\mathrm{dR}}\,d_{\mathrm{int}}=0$ by
the Leibniz rule.  The remaining five terms pair into three
cancellation mechanisms, governed by the combinatorics of index
overlap:
\begin{enumerate}[label=(\alph*)]
\item \emph{Disjoint supports.}\; When $\{i,j\}\cap\{k,l\}
=\emptyset$, the boundary faces~$D_{\{i,j\}}$
and~$D_{\{k,l\}}$ are transverse in
$\overline{\mathrm{FM}}_n(X)$.  The corresponding residue
operators commute, so the cross-terms cancel:
$d_{\mathrm{res}}^{(ij)}\,d_{\mathrm{res}}^{(kl)}
+d_{\mathrm{res}}^{(kl)}\,d_{\mathrm{res}}^{(ij)}=0$.
Operadic associativity of disjoint contractions.
\item \emph{Shared index.}\; When
$\{i,j\}\cap\{k,l\}=\{j=k\}$, the triple collision
$z_i=z_j=z_l$ lies in a codimension-two stratum.  The
form-theoretic contribution involves
$\eta_{ij}\wedge\eta_{jl}$.  The Arnold relation
\[
\eta_{ij}\wedge\eta_{jl}
+\eta_{jl}\wedge\eta_{li}
+\eta_{li}\wedge\eta_{ij}
\;=\;0
\]
identifies this with a cyclic sum over the three ordered
approaches to the triple-collision point.  On the algebraic
side, the three terms produce
$(a_{i(m)}a_j)_{(n)}a_l$, $(a_{j(m)}a_l)_{(n)}a_i$,
$(a_{l(m)}a_i)_{(n)}a_j$ (with appropriate pole orders),
whose sum vanishes by the Borcherds identity.  The Arnold relation kills the form; the Borcherds identity
kills the algebra.  Both are needed, and both are consumed.
\item \emph{Same pair.}\;
$d_{\mathrm{res}}^{(ij)}\,d_{\mathrm{res}}^{(ij)}=0$ by the
sign rule on desuspensions: extracting a residue twice at the
same diagonal is zero for degree reasons.
\end{enumerate}

The result: $d_{\barB}^2=0$ if and only if the $n$-th products
satisfy the Borcherds identity.  The bar construction accepts
exactly chiral algebras.

All of this is genus~$0$.  At genus~$g \geq 1$, the propagator acquires
quasi-periodicity corrections and $d_{\mathrm{fib}}^{\,2} \neq 0$:
the curvature $d_{\mathrm{fib}}^{\,2} = \kappa(\cA) \cdot \omega_g$
produces a single scalar $\kappa(\cA)$ controlling the genus
expansion.  Genus~$0$ is uncurved algebra;
genus~$\geq 1$ is curved geometry.


% ====================================================================
\subsection*{The categorical logarithm and two duals}
% ====================================================================

The bar construction is a categorical logarithm: it sends tensor
products to direct sums.
\[
B(\cA\otimes\cA')
\;\simeq\;
B(\cA)\oplus B(\cA').
\]
The logarithmic form $\eta_{ij}=d\log(z_i-z_j)$ is the integral
kernel.  The Arnold relation is the cocycle condition for~$\log$:
the Arnold relations make $d\log$ well-defined on products.

The cobar functor $\Omega(\barB(\cA))$ is the exponential.  The
domain on which the exponential inverts the logarithm, the
\emph{radius of convergence}, is the Koszul locus: the full
subcategory of chiral algebras for which
\[
\varepsilon\colon\Omega(\barB(\cA))
\;\xrightarrow{\;\sim\;}\;
\cA.
\]
Outside the Koszul locus, the cobar-bar resolution is infinite:
higher homotopies measure the failure of convergence.

The cobar functor recovers the original algebra~$\cA$ from its
bar complex.  A different operation, linear duality of bar cohomology,
produces a different algebra: the Koszul dual~$\cA^!$.

The \emph{Koszul dual algebra}~$\cA^!$ is obtained by a two-step
algebraic procedure.  First, take bar cohomology:
\[
\cA^{\mathrm i}
\;:=\;
H^*\bigl(B(\cA)\bigr),
\]
a graded coalgebra (the coproduct of $B(\cA)$ descends to
cohomology).  Second, apply graded linear duality:
\[
\cA^!
\;:=\;
\bigl(\cA^{\mathrm i}\bigr)^\vee.
\]
The result is the Koszul dual algebra, not the cobar.
$\Omega(B(\cA))\simeq\cA$ is bar-cobar \emph{inversion},
recovering the original algebra; $\cA^!$ is obtained by
\emph{linear duality of bar cohomology}, producing a different
algebra except when $\cA\simeq\cA^!$ (Koszul self-duality; e.g.\ $\mathrm{Vir}_{c=13}$).

A separate, geometric operation also acts on the bar coalgebra:
Verdier duality on $\mathrm{Ran}(X)$ sends $\barB_X(\cA)$ to the
\emph{homotopy} Koszul dual~$\cA^!_\infty$:
\[
\mathbb D_{\mathrm{Ran}}\,\barB_X(\cA)
\;\simeq\;
\cA^!_\infty.
\]
Verdier duality converts the coalgebra structure of~$\barB_X(\cA)$
into an algebra structure, so $\cA^!_\infty$ is a factorization
\emph{algebra}.  It is richer than~$\cA^!$: it retains full
chain-level data, while~$\cA^!$ remembers only cohomology.  That
the cohomology of $\cA^!_\infty$ recovers~$\cA^!$ on the Koszul
locus is the content of Theorem~A (the Verdier intertwining
theorem), not a tautology.

When~$\cA$ is chirally Koszul, $B(\cA)$ has cohomology concentrated
in a single row, $\cA^{\mathrm i}$ is a coalgebra with a unique
grading, and $\cA^!$ governs a linear (quadratic) resolution
of~$\cA$.  The bar-to-cobar spectral sequence collapses at~$E_2$.
Classical Koszul duality
(Priddy, Beilinson--Ginzburg--Soergel, Loday--Vallette) is
recovered on $X=\mathbb A^1$: the affine line is contractible,
so $\overline{C}_n(\mathbb A^1)$ deformation-retracts onto
a point, the FM chain cooperad becomes quasi-isomorphic to the
associative cooperad, and the bar construction reduces to the
classical bar complex of the underlying Lie algebra.


% ====================================================================
\subsection*{Two Koszul dualities: the ordered bar as primitive}
% ====================================================================

The bar complex of a chiral algebra~$\cA$ carries two operadic
colours, and each colour has its own Koszul duality.

\medskip
\noindent\textbf{The ordered bar (primitive).}\enspace
The classical bar construction of Stasheff assigns to any
$A_\infty$-algebra the cofree \emph{tensor} coalgebra
$\barB^{\mathrm{ord}}(\cA) = T^c(s^{-1}\bar\cA)$ equipped with the
deconcatenation coproduct.
This is an $E_1$-coalgebra: the coproduct is coassociative but
\emph{not} cocommutative, reflecting the linear ordering of the
tensor factors.  It is the natural primitive.  On the chiral side,
$\barB^{\mathrm{ord}}(\cA)$ is built on ordered configurations
$\operatorname{Conf}_n^{<}(\mathbb R)$ parametrising the positions of
line operators on an interval; the linear ordering on~$\mathbb R$ is
retained, with no $\Sigma_n$-quotient.  Its Koszul
dual~$\cA^!_{\mathrm{line}}$ governs the category of line operators:
when~$\cA$ is affine Kac--Moody, $\cA^!_{\mathrm{line}}$ is the
Yangian~$Y(\mathfrak g)$ (the quantum group that deforms the
enveloping algebra $U(\mathfrak g[z])$ of polynomial currents).

\medskip
\noindent\textbf{The symmetric bar (coinvariant quotient).}\enspace
The \emph{symmetric} bar complex
$\barB^{\Sigma}(\cA) = \operatorname{Sym}^c(s^{-1}\bar\cA)$ is the
$\Sigma_n$-coinvariant quotient of the ordered bar.  It carries the
\emph{coshuffle} coproduct: at tensor degree~$n$, the coshuffle sums
over all $2^n$ bipartitions of an \emph{unordered} set, in contrast to
the $n+1$ consecutive splittings of the deconcatenation.  This is an
$E_\infty$-coalgebra: coassociative and cocommutative.  Geometrically,
it is built on the Fulton--MacPherson compactification
$\overline{\operatorname{FM}}_n(\mathbb C)$ with
$\Sigma_n$-coinvariants, and it is the natural coalgebra on
$\operatorname{Ran}(X)$: the factorization coalgebra of Beilinson and
Drinfeld.  Its Koszul dual~$\cA^!_{\mathrm{ch}}$ is the chiral Koszul
dual described above, the Francis--Gaitsgory construction via Verdier
duality on~$\operatorname{Ran}(X)$.

\medskip
\noindent\textbf{The $R$-matrix as cross-colour data.}\enspace
The averaging map
$\mathrm{av}\colon\barB^{\mathrm{ord}}(\cA)\to\barB^{\Sigma}(\cA)$
takes $\Sigma_n$-coinvariants.  For a chiral algebra with OPE poles,
this descent requires the spectral $R$-matrix
$R(z)$ as twisting datum.  The $R$-matrix is the arity-$2$ component
of the ordered Maurer--Cartan element~$\Theta^{E_1}_{\cA}$: it is
the collision residue
\[
r(z) \;=\; \operatorname{Res}^{\mathrm{coll}}_{0,2}
\bigl(\Theta^{E_1}_{\cA}\bigr),
\]
a matrix-valued rational function of the spectral parameter~$z$
encoding the full binary OPE data.  Applying the averaging map at
arity~$2$ collapses the $z$-dependent profile to a single scalar:
\[
\mathrm{av}\bigl(r(z)\bigr) \;=\; \kappa(\cA),
\]
the modular characteristic of Theorem~D\@.  More generally, each
arity-$n$ component of the ordered Maurer--Cartan element projects
under~$\mathrm{av}$ to the arity-$n$ shadow of the modular
obstruction tower~$\Theta_{\cA}$: at arity~$3$, the KZ associator
projects to the cubic shadow~$C$; at arity~$4$, the Yangian higher
coproduct projects to the quartic resonance class~$Q$.

The five main theorems~A--D and~H are the invariants of the ordered
Maurer--Cartan element~$\Theta^{E_1}_{\cA}$ that survive averaging.
The modular characteristic~$\kappa$, the shadow obstruction tower, the
complementarity pairing, and the chiral Hochschild cohomology are all
$\Sigma_n$-coinvariant projections of the richer $E_1$-data.  The
ordered story is the primitive; the modular story is its shadow.

The $R$-matrix is derived from the local OPE for vertex algebras
($E_\infty$-chiral), and is independent input for nonlocal quantum
vertex algebras ($E_1$-chiral, in the sense of Etingof--Kazhdan).

% ====================================================================
\subsection*{The Heisenberg atom}
% ====================================================================

The Heisenberg chiral algebra $\mathcal H_k$ is generated by a
single field~$\alpha$ of conformal weight~$1$ with OPE
\[
\alpha(z)\,\alpha(w)
\;\sim\;
\frac{k}{(z-w)^2}\,.
\]
Only the first product is nonzero: $\alpha_{(1)}\alpha=k\cdot\mathbf 1$
and $\alpha_{(n)}\alpha=0$ for $n\neq 1$.  No first-order pole:
$\alpha_{(0)}\alpha=0$.  The Borcherds identity is satisfied
trivially: all triple products vanish since the only nontrivial
product is $\alpha_{(1)}\alpha\in\Bbbk\cdot\mathbf 1$, and
$\mathbf 1_{(n)}=0$ for $n\ge 0$.

The bar complex at tensor degree~$n$:
\[
\barB^n_X(\mathcal H_k)
\;=\;
(s^{-1}\alpha)^{\otimes n}
\;\otimes\;
\Gamma\bigl(\overline{\mathrm{FM}}_n(X),\;
\Omega^{n-1}_{\log}\bigr).
\]
\emph{Tensor degree~$1$}: the space is $s^{-1}\alpha$, one-dimensional,
with $d_{\barB}=d_{\mathrm{int}}=0$.

\emph{Tensor degree~$2$}: the bar component is
$(s^{-1}\alpha)^{\otimes 2}\otimes\langle\eta_{12}\rangle$.
The bar differential extracts the full chiral product:
$d_{\mathrm{res}}(s^{-1}\alpha\otimes s^{-1}\alpha
\otimes\eta_{12}) = \alpha_{(1)}\alpha = k\cdot\mathbf{1}$.
The double-pole mode is extracted by
$d_{\mathrm{curvature}}$
(Proposition~\ref{prop:pole-decomposition}); since
$\alpha_{(0)}\alpha = 0$, the bracket component vanishes.
The level~$k$ is visible already at genus~$0$; at genus~$1$
the curvature acquires a topological component
$d_{\mathrm{fib}}^2 = \kappa(\cH_k)\cdot\omega_1$
(Example~\ref{ex:heisenberg-d-deg1}).

\emph{Tensor degree~$n\ge 3$}: every residue at a codimension-one
collision stratum $D_{\{i,j\}}$ extracts $\alpha_{(1)}\alpha
=k\cdot\mathbf 1$ on the pair $(i,j)$.  The remaining $n-2$
factors are~$\alpha$; since $\mathbf 1_{(p)}\alpha=0$ for
$p\ge 0$, no further contractions occur.  All structure is
captured at tensor degree~$2$: the bar complex is
\emph{quadratic}.  The Heisenberg algebra is chirally Koszul.

The modular characteristic:
\[
\kappa(\mathcal H_k)\;=\;k.
\]
The Koszul dual is the \emph{curved} symmetric chiral algebra
\[
\mathcal H_k^!
\;=\;
\bigl(\mathrm{Sym}^{\mathrm{ch}}(V^*),\;m_0=-k\omega\bigr),
\]
where $V^*$ is the $\Bbbk$-linear dual of the generating space and
$m_0=-k\omega$ is the curvature: a degree-$2$ element satisfying
the curved $A_\infty$ relation
\[
m_1^{\,2}(a)\;=\;[m_0,a]
\]
(commutator, with the minus sign).  The differential $m_1$ does not
square to zero; its failure is controlled by the curvature.  The
curvature~$m_0$ is the bar-dual of the OPE coefficient~$k$:
the algebra side has a scalar OPE; the dual side has a scalar
curvature.

The Heisenberg has only a second-order pole: the bar complex is
quadratic.  What Chriss--Ginzburg demands at this point is the simplest
failure: the first example where the bar complex is \emph{not} quadratic,
where the first-order pole generates genuinely new structure.

\medskip
\noindent\textbf{The $E_1$ face.}\enspace
The ordered bar complex $\barB^{\mathrm{ord}}(\mathcal H_k)$ carries
the $R$-matrix
\[
r(z) \;=\; \frac{k}{z}\,,
\]
a scalar rational function with a single simple pole at the origin.
This is the collision residue of the ordered Maurer--Cartan element at
arity~$2$: the second-order OPE pole $k/(z-w)^2$ drops by one order
under the logarithmic kernel $d\log(z-w)$ (the bar differential
absorbs one power of $(z-w)$), producing the first-order $r$-matrix
pole $k/z$.  Applying the averaging map:
\[
\mathrm{av}\bigl(r(z)\bigr) \;=\; k \;=\; \kappa(\mathcal H_k).
\]
For the Heisenberg algebra, the $R$-matrix \emph{is} the modular
characteristic: averaging loses nothing, because $r(z) = k/z$ is
already $\Sigma_2$-invariant (a scalar, not a matrix).  This is the
CG opening: the simplest example where the two Koszul dualities
coincide, the ordered bar carries no more information than the
symmetric bar, and the entire structure collapses to a single number~$k$.
The deficiency is immediate: no first-order pole, no Lie bracket, no
quantum group.


% ====================================================================
\subsection*{The Kac--Moody OPE}
% ====================================================================

The affine Kac--Moody algebra $\widehat{\mathfrak g}_k$ at level~$k$
is generated by currents $J^a$, $a=1,\dots,\dim\mathfrak g$, with
OPE
\[
J^a(z)\,J^b(w)
\;\sim\;
\frac{f^{ab}{}_{c}\,J^c(w)}{z-w}
\;+\;
\frac{k\,\kappa^{ab}}{(z-w)^2}\,,
\]
where $f^{ab}{}_{c}$ are the structure constants of the simple Lie
algebra~$\mathfrak g$ and $\kappa^{ab}$ is the Killing form
normalised so that long roots have squared length~$2$.  Both poles
are present: the first-order pole carries the Lie bracket
$J^a_{(0)}J^b=f^{ab}{}_{c}J^c$; the second-order pole carries
the invariant inner product
$J^a_{(1)}J^b=k\,\kappa^{ab}\cdot\mathbf 1$.  The bar complex
of $\widehat{\mathfrak g}_k$ has nontrivial contributions
at all tensor degrees: the first-order pole generates structure
constants at tensor degree~$n$ from tree-level OPE contractions,
and the second-order pole contributes the metric/curvature data.

The bracket and curvature decomposition:
\[
\mu^{\mathrm{ch}}
\;=\;
\underbrace{f^{ab}{}_{c}\,J^c\cdot\eta_{12}}_{\text{Lie bracket
(first-order pole)}}
\;+\;
\underbrace{k\,\kappa^{ab}\cdot\mathbf 1\cdot
d\eta_{12}}_{\text{inner product (second-order pole)}}.
\]
The Lie bracket contributes the tree-level bar differential
(combinatorics of the Chevalley--Eilenberg complex); the inner
product contributes the curvature~$\kappa$.  The coexistence of
both poles, absent in the pure second-order-pole Heisenberg case,
makes $\widehat{\mathfrak g}_k$ the first nontrivial example.

The ordered bar complex of $\widehat{\mathfrak g}_k$ carries the $R$-matrix
\[
r(z) \;=\; \frac{\Omega}{z}\,,
\qquad
\Omega \;=\; \sum_a J^a \otimes J_a\,,
\]
where $\Omega$ is the Casimir tensor of~$\mathfrak g$.  This is
matrix-valued: $\Omega \in \mathfrak g \otimes \mathfrak g$
carries the full Lie-algebraic structure.  Averaging collapses the
Casimir to its trace:
\[
\mathrm{av}\bigl(\Omega/z\bigr)
\;=\;
\frac{\dim(\mathfrak g)\cdot(k+h^\vee)}{2h^\vee}
\;=\;
\kappa(\widehat{\mathfrak g}_k).
\]
The Casimir tensor structure, invisible to the modular characteristic,
is the data that builds the Yangian $Y(\mathfrak g)$.  At this point
the ordered bar has strictly more information than the symmetric bar:
the passage from $r(z)=\Omega/z$ to $\kappa(\widehat{\mathfrak g}_k)$
discards the matrix structure of~$\Omega$, which encodes the representation
theory of~$\mathfrak g$.

\begin{remark}[The zero-dimensional shadow]\label{rem:pf1-zero-dim}
When $X=\operatorname{Spec}\Bbbk$ is a point, configuration spaces
collapse, logarithmic forms vanish, and the chiral bar complex
recovers the classical bar construction of an associative
algebra~$A$.  A \emph{quadratic algebra} $A=T(V)/(R)$ has quadratic
dual $A^!=T(V^*)/(R^\perp)$; when~$A$ is Koszul, the bar
cohomology concentrates in diagonal degrees and
$H^*(B(A))\cong(A^!)^\vee$.  The operadic instances
$\mathrm{Com}^!=\mathrm{Lie}$ and $\mathrm{Sym}^!=\Lambda$ recover
the Chevalley--Eilenberg complex as a restricted bar construction.
When Koszulness fails, the cobar-bar resolution
$\Omega(B(A))\xrightarrow{\sim}A$ is the universal replacement.
\end{remark}

Everything above takes place at genus~$0$: the propagator
$\eta_{12}=d\log(z_1-z_2)$ is globally defined on~$\mathbb P^1$,
and the Arnold relation holds without correction.  On a curve of
genus~$g \geq 1$, the function $z_1 - z_2$ has no global meaning:
the curve is compact, and any meromorphic function with a single
simple zero must also have a simple pole.  The Arnold relation
acquires a defect, the bar differential no longer squares to zero,
and the defect is controlled by a single scalar.
